\chapter{Introduction}

\section{Background and Motivation}

Cloud computing adoption has accelerated significantly with major vendors such as Amazon Web Services (AWS) and Microsoft Azure continually expanding proprietary service offerings. These cloud platforms provide a wide range of infrastructure, platform, and software services catering to diverse business needs. However, the growing variety and complexity of cloud offerings introduce significant challenges for solution architects who must balance competing factors such as cost, scalability, security, and performance when designing cloud architectures. %\cite{agiliway2025aws-vs-azure}.

The present study proposes an innovative agentic AI framework to assist architects in navigating the resulting architectural complexity. The framework employs an Evaluator-Optimizer workflow \cite{anthropic2024agents} that iteratively generates, validates, and refines cloud architectural recommendations. Validation is grounded in authoritative cloud provider documentation through Retrieval Augmented Generation (RAG) \cite{lewis2020rag, asai2023selfrag}, ensuring factual and reliable outputs. The system aims to deliver efficient, secure, and cost-optimized solutions capable of cross-provider comparative analysis. %\cite{futransolutions2025multi-cloud}.

\section{Problem Statement}

The transition from traditional on-premises infrastructure to cloud platforms has resulted in an exponential increase in cloud services. Providers like AWS and Azure often offer overlapping services solving similar problems, complicating architectural decisions. %\cite{simplilearn2025.aws-vs-azure-challenges}
Architects typically rely on domain expertise and experience to evaluate tradeoffs involving costs, scalability, security, and performance. Current AI-based tools and assistants, including first-party vendor AI products like Amazon Q \cite{aws2023q} and Microsoft Copilot \cite{microsoft2024copilot}, suffer from several limitations:

\begin{itemize}
    \item Existing chatbots often provide partial guidance focusing on isolated components rather than comprehensive end-to-end architectures.
    \item Generated outputs frequently lack grounded factual validation against authoritative cloud documentation---and retrieval quality degrades sharply when dense technical terminology is involved, an issue that naive vector search alone cannot resolve.
    \item Current solutions generally do not support comparative evaluation across cloud providers, nor do they offer structured mechanisms to surface provider-specific tradeoffs for informed decision-making.
    \item Without a coherent user interface, solution architects cannot inspect the reasoning process, switch between provider recommendations, or iterate with visual feedback during a session.
\end{itemize}

The outlined challenges highlight the need for an intelligent AI-driven framework capable of generating, validating, and optimizing architectural recommendations with multi-cloud support, advanced retrieval quality, and an LLM-as-a-judge evaluation protocol for rigorous assessment. %\cite{agiliway2025aws-vs-azure, futransolutions2025multi-cloud}.

\section{Objectives}

The primary objectives of this research project are:

\begin{itemize}
    \item To design and implement an Evaluator-Optimizer feedback workflow using advanced Retrieval Augmented Generation (RAG) enhanced with cross-encoder reranking for factually grounded recommendation validation.
    \item To extend the system to a full multi-cloud architecture supporting both AWS and Azure through a provider metadata registry that cleanly separates vendor-specific configuration from the core agent logic.
    \item To incorporate cost and feature tradeoff analysis to enhance efficiency, security, and cost optimization in architectural decisions.
    \item To expose Infrastructure-as-Code generation hints (Terraform, CloudFormation, Bicep) as a first step towards directly deployable output artifacts.
    \item To evaluate framework quality using an LLM-as-a-judge protocol across six technical dimensions, replacing the surface-level lexical similarity metrics used in Phase 1.
    \item To develop a responsive Next.js user interface providing real-time workflow visualization, side-by-side provider comparison, and per-run model selection across the reasoning and execution tiers.
\end{itemize}

\section{Scope and Limitations}

\subsection{Scope}

The research focuses on developing and evaluating a functional prototype of the agentic AI framework, limited to the following boundaries:

\begin{itemize}
    \item \textbf{Cloud Providers:} Initially limited to Amazon Web Services (AWS) and Microsoft Azure, the two dominant public cloud vendors.
    \item \textbf{Knowledge Source:} The Retrieval Augmented Generation knowledge base is constructed exclusively from publicly available, official provider documentation such as whitepapers, service guides, and recommended best practices.
    \item \textbf{Output Format:} The system outputs comprehensive textual cloud-architecture recommendations with detailed justifications, as well as generates deployable code (Infrastructure-as-Code templates such as Terraform or CloudFormation).
\end{itemize}

\subsection{Limitations}

The implemented prototype currently has the following operational constraints:

\begin{itemize}
    \item \textbf{Absence of Real-Time Data:} The framework does not incorporate real-time information such as pricing updates, service availability in specific regions, or customer-specific account configurations.
    \item \textbf{Knowledge Base Boundaries:} Although extensive, the knowledge base may lack coverage of some niche or recently introduced cloud services due to document availability and constraints associated with the open source vector database.
    \item \textbf{Non-Substitutive Advisory Role:} The system is intended as a decision support tool complementing human expertise rather than replacing human operators. All architectural recommendations warrant expert review prior to deployment.
\end{itemize}

\section{Feasibility Study}

\subsection{Technical Feasibility}

Mature underlying technologies such as Large Language Models (LLMs), vector similarity searches, and agentic AI frameworks (e.g., LangChain \cite{langchain2024docs}, LangGraph \cite{langgraph2024docs}, AutoGen \cite{wu2023autogen}) enable the realization of the proposed framework. Cloud platforms (AWS and Azure) provide rich documentation and resources essential for building the knowledge base.

\subsection{Data Availability}

Cloud vendor documentation required for the RAG knowledge base is openly accessible. Such resources include service whitepapers, architecture best practices, and official guides, forming a comprehensive foundation for reliable information retrieval and validation. %\cite{aws2025well-architected, azure2025-architecture}.

\subsection{Operational Feasibility}

The prototype targets deployment on affordable local computation platforms using GPU-accelerated frameworks (Apple Metal Performance Shaders \cite{apple2024mps}) and local embeddings (Ollama \cite{ollama2024web}), with the potential to integrate cloud based services. The user interfaces will be minimal yet practical comprising command-line and web portals, to streamline accessibility, including for evaluation purposes.

\subsection{Expertise and Challenges}

The project benefits from the author's background in data engineering, AI/ML, cloud computing, and Python development. The primary anticipated challenge is ensuring high accuracy and factual reliability of AI-generated recommendations, requiring robust knowledge ingestion and validation techniques.

\section{Thesis Organization}

The dissertation is organized into the following chapters:

\begin{itemize}
    \item \textbf{Chapter 2} provides a literature review covering foundational concepts: agentic multi-agent systems, Retrieval-Augmented Generation, advanced retrieval techniques including cross-encoder reranking, multi-agent debate paradigms, and LLM-as-a-judge evaluation methodology.
    \item \textbf{Chapter 3} details the system design and methodology, describing the subgraph-based LangGraph orchestration, the provider metadata registry enabling multi-cloud extensibility, and the knowledge base construction pipeline for both AWS and Azure.
    \item \textbf{Chapter 4} narrates the implementation of all major components: vector store construction with FlashRank reranking, the modular provider metadata and IaC hint system, the State TypedDict reducer design, the dynamic model selection mechanism, the Next.js user interface, and the evaluation infrastructure comprising three runner variants and five benchmark scenarios.
    \item \textbf{Chapter 5} presents empirical results: an ablation study across three framework variants (baseline, agentic draft, full Evaluator-Optimizer loop) using GPT-5.4, and a cross-model comparison of GPT-5.4, Claude Sonnet~4.6, and Gemini~3.1 Pro under the LLM-as-a-judge protocol, with discussion of methodology limitations.
    \item \textbf{Chapter 6} concludes the thesis with a summary of contributions, acknowledged limitations, and future directions including GCP provider integration and live Infrastructure-as-Code deployment.
\end{itemize}